
% ============================================================================
\clearpage
\section*{v4.10.0-r20 Source-Only: A Null Instant of the Field Invariant as a Candidate Shape-(iii) Witness for C5}
\addcontentsline{toc}{section}{v4.10.0-r20 Source-Only: A Null Instant of the Field Invariant as a Candidate Shape-(iii) Witness for C5}
\label{sec:v4100-r20-c5-null-instant}
% ============================================================================

\paragraph{Purpose.}
Layer v4.7.0-r88 left shape (iii), a witness outside the collar iteration, as the only open route to a finite C5 discharge, with
the requirements (\(\alpha\)) vanishing of the obstruction itself, not smallness of the score; (\(\beta\)) independence of
\(\mathcal U_5^{\rm bd}\); (\(\gamma\)) compatibility with \(C5_{\rm det/eff}\); (\(\delta\)) no import of the all-readout slot.  Layer
v4.9.0-r10 read the obstruction as behaving like the invariant \(s\) and concluded that a witness would need dynamics; v4.9.0-r27
proposed a kick-and-measure dynamics that gives no finite zero of the score.  This layer records a candidate from the free time
evolution of time-harmonic fields, checks it against (\(\alpha\))--(\(\delta\)), and names what it does not supply.  Exploration, not
blind.  \textbf{The forbidden assertion \(C5=0\) is not made.}  The notes are
\texttt{notes/\allowbreak research/\allowbreak v4\_10\_0\_r20/} and
\texttt{notes/\allowbreak research/\allowbreak c5\_null\_instant\_2026\_11\_15/}.  Earlier layers are not rewritten.

\noindent\textbf{Source-only status.}
This v4.10.0-r20 layer is source-only.  It does not build a PDF, does not assign a new DOI,
and does not authorize a public release.  The v5.1.0 Version DOI
\href{https://doi.org/10.5281/zenodo.24237112}{10.5281/zenodo.24237112} and the Concept DOI
\href{https://doi.org/10.5281/zenodo.19210728}{10.5281/zenodo.19210728} are unchanged by it.

\subsection*{v4.10.0-r20 results}
\label{subsec:v4100-r20-results}

\paragraph{Theoremlet v4.10.0-r20.1 (null instants).}
For a time-harmonic field with complex amplitudes the instantaneous invariant \(s(t)=\tfrac14\bigl(s_+e^{-2i\omega t}+2w+\bar s_-e^{2i\omega t}\bigr)\)
(r16.8) runs on an ellipse around \(w/2\).
\begin{enumerate}[leftmargin=2em]
\item[(a)] For \(\hat E=(1,0,0.3i)\), \(\hat B=(0.7i,1,0)\) the state is not null (\(s_+=-1\), \(s_-=1.8\), \(w=-0.4\)), yet \(s(0)=0\) exactly with
\(|z(0)|=1.41\); the null instant recurs every half period.
\item[(b)] The condition is of codimension one: for random amplitudes the ellipse misses the origin (median relative distance \(0.15\)).
\item[(c)] The null instant is a Lorentz-invariant event (\(s\) is a scalar at every instant; checked on 200 boosts).
\item[(d)] At the null instant the state of the disk of r04 is the photon point \(P\), \(b=0\); the compensator share is \(\tfrac12\), not \(0\).  Free time
evolution thus reaches \(P\) at a finite time, unlike boosts (v4.9.0-r19.1).
\item[(e)] For real in-phase amplitudes the null instant is a zero of the whole field and is no witness.
\end{enumerate}

\paragraph{Remark v4.10.0-r20.2 (the requirements of r88; reading).}
(\(\beta\)): the free evolution is independent of the collar recursion.  (\(\gamma\)): no conflict with the detector statement was found.
(\(\delta\)): one field, one readout.  (\(\alpha\)): Reading~B considers real fields at a point (v4.9.0-r05.2(a)), and r10.2 grounds the
obstruction in the floor \(|s|\) of that instantaneous field; within the r10 reading the relevant quantity is the instantaneous \(s(t)\),
not the state invariants \(s_\pm\), \(w\).  At the null instant this floor is exactly \(0\), while content and score stay positive.  Within
that reading (\(\alpha\)) is met; the reading itself remains an interpretation.

\paragraph{Theoremlet v4.10.0-r20.3 (time and collar stage).}
A collar stage is a duality rotation followed by compensation (v4.9.0-r10.1, r27); duality turns the two helicity parts in
opposite senses, time translation in the same sense.
\begin{enumerate}[leftmargin=2em]
\item[(a)] They coincide only for pure helicity: with \(F_-=0\), \(z(t)=e^{-i\omega t}z(0)\).
\item[(b)] Pure-helicity fields have \(|s(t)|\) constant and no null instant.
\item[(c)] Fields with a null instant carry both helicities, and their time evolution is no sequence of duality rotations (relative
remainder against the best rotation, median \(0.49\)).
\end{enumerate}
Reading: the null instant arises exactly where the time evolution leaves the collar dynamics, which is the independence (\(\beta\))
asks for; it also leaves open how an instant would be counted as a finite stage \(N\) in the sense of v4.7.0-r86.

\paragraph{Theoremlet v4.10.0-r20.4 (a frame audit of C5).}
In the reading of v4.9.0-r10, r27 and r28 the C5 score is the content of the compensator port \(C_{\rm HC}=\tfrac12(1+\sigma_x)\) of the
doublet, on the axis of the side \(\operatorname{sgn}(E\cdot B)\); by r28 the tachyon map \(z\mapsto iz\) exchanges these ports and \(J_B\) fixes them.
\begin{enumerate}[leftmargin=2em]
\item[(a)] Two notions of port are in use.  The ports of the C5 reading are exchanged by \(R\) and \(T\) and fixed by \(J_B\); the ports
\(\tfrac12(1\pm iL_{e_7})\) of v4.9.0-r42 and r51, which by r16.2(b) are the \(\sigma_y\) poles, are fixed by \(R\) and exchanged by \(T\) and \(J_B\).
\item[(b)] Per frame, the selective collar recursion gives the matter frame the share \((3/4)^n\) and the tachyon frame, which reads the
exchanged ports, the share \(1-(3/4)^n\); their sum is \(1\) at every stage.  The old question \(C5=0\) is, per frame, the question whether the
tachyon frame reads the whole content: a complete change of frame, finite only by the quarter turn, which C4, C6 and (\(\delta\)) block.
\item[(c)] At a null instant (r20.1) both shares are \(\tfrac12\); in the example the matter share rises through the instant from \(0.29\) to
\(0.71\).  No frame score vanishes; the floor \(|s|\), which the tachyon map sends to itself, does.
\end{enumerate}
Reading: like OP4 before v4.9.0-r61, C5 in its old form asks, in the measure of one frame, about a passage between two frames.  A
frame-clean statement keeps the shares \(C5_M\) and \(C5_T\) with fixed sum; its frame-free point is the photon point \(P\), which free time
evolution reaches at a finite time.  Whether C5 is restated in this way is the author's decision; \(C5=0\) is not asserted.

\paragraph{What this does not give.}
No C5 discharge.  Open: a periodic rather than lasting zero (r86 asks for a finite stage \(N\) with \(\Phi_5^{(N)}=0\) and does not say
whether the zero must persist); the count of an instant as a stage; the content does not fall to zero at the null instant, and a boost
could empty it there only in the limit (r10.2(c)); why non-generic fields of this kind should be present.

\subsection*{v4.10.0-r20 status in the three frames}
\label{subsec:v4100-r20-frames}

\begin{itemize}[leftmargin=2em]
\item M, T: the null instant is a Lorentz-invariant event; the tachyon map turns \(s\) into \(-s\) and keeps \(s=0\).
\item P: the state at the null instant is the photon point \(P\) (r20.1(d)).
\end{itemize}

\subsection*{v4.10.0-r20 register entries}
\label{subsec:v4100-r20-register}

\begin{center}
\small
\begin{tabular}{@{}>{\raggedright\arraybackslash}p{0.34\textwidth}>{\raggedright\arraybackslash}p{0.58\textwidth}@{}}
\toprule
statement & status after v4.10.0-r20\\
\midrule
null instant of \(s(t)\) & exists for a codimension-one family; invariant; state at \(P\), score \(\tfrac12\) (r20.1)\\
C5 shape (iii), requirements of r88 & (\(\beta\))--(\(\delta\)) met; (\(\alpha\)) within the r10 reading (r20.2)\\
time and collar stage & coincide only for pure helicity, which has no null instant (r20.3)\\
frame audit of C5 & old \(C5=0\) is a complete change of frame; frame-free point \(P\); two notions of port (r20.4)\\
\(C5_{\rm zero}\) & open; no discharge, no assertion \(C5=0\)\\
\bottomrule
\end{tabular}
\end{center}

\subsection*{v4.10.0-r20 ledger and guard update}
\label{subsec:v4100-r20-ledger-guard}

\[
R_{\rm chk}^{v4100\text{-}r20}=(1,1,0)
\]
for (candidate recorded; computations checked in the cited scripts; persistence, stage count and the presence of such fields open).

The layer spends no guard:
\[
R_{\rm glob}^{(v4100\text{-}r20,+)}=(0,1,1,1,1),
\qquad
\Gamma_{\rm guard}^{v4100\text{-}r20}=(1,1,1,1,1),
\qquad
\beta_4=1/256\ \text{unspent}.
\]
The layer does not imply exact finite \(C5\) zero closure, exact all-readout terminality, \(C6/T2\), a global photon-clock/RH theorem,
or global all-readout closure.
